\begin{section}{Introduction to Topic and Research Question}
\begin{subsection}{Background and Motivation}
    Machine Learning (ML) has evolved into a central paradigm in computer science, providing methods to
    automatically detect meaningful patterns in data and and make predictions based on these patterns.  In contrast
    to traditional programmed algorithms, learning systems are capable of adapting to new data and change 
    their behavior accordingly\footcite[Cf.][p. 128]{osisanwo2017}. This data-driven approach allows computers to solve problems that are too complex
    to tackle with handcrafted methods, and it underpins real-world applications. For instance, \enquote{Smart classifiers protect our email by learning from
    massive amounts of spam data and user feedback; advertising systems learn to match the right ads with the right context; fraud detection systems 
    protect banks from malicious attackers; anomaly event detection systems help experimental physicists to find events that lead to new physics.} \footcite{chen2016}
    - tasks that rely on extracting insights from large datasets rather than explicit programming.

    Among the fundamental approaches in ML is \textit{supervised learning}, were a model is trained on labeled
    examples (input-output pairs) to predict the correct output for new, unseen inputs. The learning aims 
    to generalize from the training data and make accurate predictions on future data\footcite[Cf.][p. 249]{kotsiantis2007}.
    A common use case is \textbf{classification}, especially binary classification, where the algorithm assigns
    each instance to one of two possible categories (e.g. ``postitive'' vs. ``negative'') based on its features\footcite[Cf.][p. 249]{kotsiantis2007}.
    Binary classification is a basic yet critical task in ML, forming the basis for many decision making systems. 
    Given its prevalence, improving the accuracy and efficiency of binary classification models is of high importance.
\end{subsection}
\begin{subsection}{Problem Statement}
    In modern machine learning practice, it is common to combine multiple models to improve predictive
    performance. Such techniques are known as \textit{ensemble methods}, which aggregate the outputs of several base
    learners to achieve better accuracy and robustness than a single model alone\footcite[Cf.][pp. 1--2]{natekin2013}.
    One successful ensemble method is \textbf{boosting}, where models are added iteratively, so that each
    new model can correct the errors of the ensemble built so far \footcite[Cf.][pp. 1--2]{natekin2013}.
    Gradient boosting methods have demonstrated state-of-the-art results on many classification tasks by iteratively fitting new 
    models to the residual errors of the previous models.

    \textbf{XGBoost} (eXtreme Gradient Boosting) is a advanced implementation of gradient boosting that introduces
    several optimizations to enhance both predictive accuracy and computational efficiency. In addition to the
    standard boosting procedure, XGBoost includes innovations like explicit regularazation of model
    complexity, parallelized tree construction, and cache-aware data structures for faster access, which makes it 
    a highly efficient and powerful ensemble learner\footcite[Cf.][pp. 785--786]{chen2016}.
    Within the XGBoost framework, the choice of \textit{booster} (the type of base learner) significantly influences the 
    model's behavior and performance. The library offers two primary booster options: \textbf{gbtree}, which uses decision trees
    capable of capturing complex, nonlinear relationships in the data. The linear booster \textbf{gblinear}, produces
    an additive linear model (with weight regularazation) analogous to logistic or linear regression \footcite[Cf.]{xgboost2024}.
    In other words, \code{gbtree} booster relies on treebased models, while \code{gblinear} uses linear functions as base learners.\footcite[Cf. ]{xgboost2024}.

    Tree-based ensembles are generally known for their strong predictive power across diverse tasks \footcite[Cf.][p. 785]{chen2016}, whereas 
    linear models tend to train faster and can perform well in high-dimensional feature spaces (for example, text classification with thousands of sparse features) \footcite[Cf.][p. 129]{osisanwo2017}.
    While XGBoost's tree booster is often observed to achieve superior accuracy on complex datasets, the linear booster
    may has advantages in terms of training speed and simplicity, especially when the data relationships are linear.
    It remains an open question how these two booster types compare head-to-head in a sprecific domain. Therefore, This
    study aims to comparatively evaluate \textbf{gbtree} and \textbf{gblinear} in terms of their predictive accuracy and efficiency
    on a binary user classification task. By analyzing their performance under the same conditions, the work seeks to identify the trade-offer
    between the nonlinear modeling capacity of tree ensembles and the potential speed or simplicity benefits of linear models. 
    This investigation will inform practitioners about which booster is better suited for a given problem setting
    and contribute to a deeper understanding of XGBoost's capabilities.
\end{subsection}
\end{section}
